\chapter{Point Splitting and Stress Tensor Renormalization}
\label{ch:point-splitting-stress-tensor-renormalization}

\ConventionsMixed

The previous chapter introduced Hadamard states as the admissible
short-distance class for free fields on curved backgrounds.  Point splitting
uses the state-independent singular part of a Hadamard two-point function to
define composite local fields.

\section{Wick Square}

\begin{definition}[Wick square by point splitting]
\label{def:wick-square-point-splitting}
Let \(\omega\) be a Hadamard state of the real Klein--Gordon field.  Let
\(H(x,y)\) denote the locally covariant Hadamard parametrix, containing the
singular \(U/\sigma_\epsilon\) and \(V\log\sigma_\epsilon\) terms.  The
point-split Wick square in state \(\omega\) is the distribution defined by
\[
  \omega\bigl(:\!\Phi^2\!:(f)\bigr)
  =
  \int_M f(x)
  \lim_{y\to x}
  \bigl(\omega_2(x,y)-H(x,y)\bigr)\,d{\rm vol}_g(x).
\]
\end{definition}

For higher Wick powers, one subtracts the parametrix in every contraction.
The result is an operator-valued distribution in perturbative algebraic QFT
and a quadratic form in representations where the required domains are
available.

The subtraction has finite local freedom.  For the Wick square in four
dimensions,
\[
  :\!\Phi^2\!:\ \mapsto\ :\!\Phi^2\!:+c_1R+c_2m^2
\]
with constants \(c_1,c_2\), subject to the chosen covariance and scaling
conditions.

\section{Stress Tensor by Bidifferential Operator}

For the classical scalar action
\[
  S[\phi]
  =
  -\frac{1}{2}\int_M
  \left(
  \nabla_\mu\phi\nabla^\mu\phi
  +m^2\phi^2+\xi R\phi^2
  \right)d{\rm vol}_g
\]
in mostly-plus Lorentzian convention, the stress tensor is obtained by metric
variation.  Quantum mechanically, define a bidifferential operator
\(\mathcal D_{\mu\nu}(x,y)\) by polarizing the classical quadratic expression
in two field arguments.  Then
\[
  \omega(T_{\mu\nu}(x))
  =
  \lim_{y\to x}
  \mathcal D_{\mu\nu}(x,y)
  \bigl(\omega_2(x,y)-H(x,y)\bigr)
  +C_{\mu\nu}(x).
\]
The local tensor \(C_{\mu\nu}\) is the finite renormalization freedom.  It is
chosen from locally covariant conserved curvature tensors, together with
terms proportional to \(m^2\) and \(m^4\) as allowed by dimension.

\section{Worked Flat-Space Computation}

We now compute a simple case explicitly.  Let
\((M,\eta)\) be four-dimensional Minkowski spacetime with
\(\eta={\rm diag}(-1,1,1,1)\), and take the massless minimally coupled free
scalar.  The vacuum two-point distribution is
\[
  W_0(x,x')
  =
  \lim_{\epsilon\downarrow0}
  \frac{1}{4\pi^2}
  \frac{1}{
  |\mathbf x-\mathbf x'|^2-(t-t'-\ii\epsilon)^2}.
\]
In flat space the local Hadamard parametrix with the same \(i\epsilon\)
prescription is \(H_0=W_0\).  Therefore the renormalized vacuum stress
tensor obtained by point splitting is zero:
\[
  \langle T_{\mu\nu}\rangle_{0,{\rm ren}}
  =
  \lim_{x'\to x}
  \mathcal D_{\mu\nu}(x,x')\bigl(W_0(x,x')-H_0(x,x')\bigr)
  =
  0 ,
\]
if the finite flat-space cosmological-constant counterterm is set to zero.
This is not an additional Lorentz-invariance argument; it is the explicit
subtraction of the state-independent short-distance singularity.

For the thermal KMS state at inverse temperature \(\beta\) with respect to
the inertial time \(t\), the difference from the vacuum two-point function is
smooth:
\[
  \Delta W_\beta(x,x')
  =
  W_\beta(x,x')-W_0(x,x')
\]
\[
  =
  \int \frac{d^3k}{(2\pi)^3}\,
  \frac{n_\beta(\omega)}{2\omega}
  \left(
  e^{-\ii\omega(t-t')+\ii\mathbf k\cdot(\mathbf x-\mathbf x')}
  +
  e^{+\ii\omega(t-t')-\ii\mathbf k\cdot(\mathbf x-\mathbf x')}
  \right),
  \qquad
  \omega=|\mathbf k|,
\]
where \(n_\beta(\omega)=(e^{\beta\omega}-1)^{-1}\).  Since
\(\Delta W_\beta\) is smooth, the point-split stress tensor is obtained by
applying the polarized classical expression directly to this remainder.
For the massless minimal tensor
\[
  T_{00}
  =
  \frac{1}{2}\bigl((\partial_t\phi)^2+|\nabla\phi|^2\bigr),
\]
the bidifferential operator is
\[
  \mathcal D_{00}
  =
  \frac{1}{2}\left(
  \partial_t\partial_{t'}
  +\sum_{i=1}^3\partial_i\partial_{i'}
  \right).
\]
On either plane wave in \(\Delta W_\beta\),
\[
  \mathcal D_{00}
  e^{-\ii\omega(t-t')+\ii\mathbf k\cdot(\mathbf x-\mathbf x')}
  =
  \omega^2
  e^{-\ii\omega(t-t')+\ii\mathbf k\cdot(\mathbf x-\mathbf x')}.
\]
Taking the diagonal limit gives
\[
  \rho_\beta
  =
  \langle T_{00}\rangle_{\beta,{\rm ren}}
  =
  \int \frac{d^3k}{(2\pi)^3}\,\omega\,n_\beta(\omega)
  =
  \frac{1}{2\pi^2}\int_0^\infty
  \frac{k^3\,dk}{e^{\beta k}-1}.
\]
The change of variables \(x=\beta k\) and the expansion
\((e^x-1)^{-1}=\sum_{n\geq1}e^{-nx}\) give
\[
  \int_0^\infty \frac{x^3\,dx}{e^x-1}
  =
  \sum_{n\geq1}\int_0^\infty x^3e^{-nx}\,dx
  =
  6\sum_{n\geq1}\frac{1}{n^4}
  =
  \frac{\pi^4}{15}.
\]
Thus
\[
  \rho_\beta=\frac{\pi^2}{30\beta^4}.
\]
For the spatial components,
\[
  T_{ij}
  =
  \partial_i\phi\,\partial_j\phi
  +\frac{1}{2}\delta_{ij}
  \bigl((\partial_t\phi)^2-|\nabla\phi|^2\bigr),
\]
so
\[
  \mathcal D_{ij}
  =
  \partial_i\partial_{j'}
  +\frac{1}{2}\delta_{ij}
  \left(
  \partial_t\partial_{t'}
  -\sum_{\ell=1}^3\partial_\ell\partial_{\ell'}
  \right).
\]
For massless modes the second term vanishes on shell because
\(\omega^2=|\mathbf k|^2\).  Rotational invariance of the momentum integral
gives
\[
  \int d\Omega\, k_i k_j=\frac{4\pi}{3}\delta_{ij}k^2,
\]
and therefore
\[
  \langle T_{ij}\rangle_{\beta,{\rm ren}}
  =
  \frac{1}{3}\rho_\beta\,\delta_{ij}.
\]
The point-split computation thus yields
\[
  \langle T_{\mu\nu}\rangle_{\beta,{\rm ren}}
  =
  \frac{\pi^2}{90\beta^4}
  \begin{pmatrix}
  3&0&0&0\\
  0&1&0&0\\
  0&0&1&0\\
  0&0&0&1
  \end{pmatrix}_{\mu\nu},
\]
in the inertial frame.  Its trace is
\(-\rho_\beta+3p_\beta=0\), as required for the massless conformal equation
of state in flat spacetime.

\section{Constant-Curvature Diagnostic}

A de Sitter-invariant Hadamard state of a free scalar has a stress-tensor
expectation value constrained by symmetry to the form
\[
  \langle T_{\mu\nu}\rangle_{\rm dS}=C\,g_{\mu\nu}.
\]
For the massless conformally coupled real scalar in four dimensions, the
state-dependent classical trace vanishes.  With the trace-anomaly convention
of Chapter~\ref{ch:curved-trace-anomalies}, the real scalar has
\(a=1/360\), \(c=1/120\), and
\[
  \langle T^\mu{}_\mu\rangle
  =
  \frac{1}{16\pi^2}\left(cW_{\mu\nu\rho\sigma}W^{\mu\nu\rho\sigma}
  -aE_4+b\Box R\right).
\]
On four-dimensional de Sitter with Hubble parameter \(H\),
\[
  W_{\mu\nu\rho\sigma}=0,\qquad
  \Box R=0,\qquad
  R=12H^2,
\]
\[
  R_{\mu\nu\rho\sigma}R^{\mu\nu\rho\sigma}=24H^4,\quad
  R_{\mu\nu}R^{\mu\nu}=36H^4,\quad
  E_4=24H^4 .
\]
Consequently
\[
  4C=\langle T^\mu{}_\mu\rangle
  =
  -\frac{H^4}{240\pi^2},
  \qquad
  C=-\frac{H^4}{960\pi^2}.
\]
This calculation isolates the part fixed by local covariance, de Sitter
invariance, and the anomaly normalization.  Massive fields, and schemes that
allow additional dimension-four metric counterterms, shift \(C\) by finite
local terms; the point-splitting prescription must state that choice.

\section{Conservation and Trace}

\begin{definition}[Conserved stress tensor]
\label{def:conserved-renormalized-stress-tensor-curved}
Within a locally covariant point-splitting scheme, a renormalized stress
tensor is called conserved when its finite local curvature terms have been
chosen so that
\[
  \nabla^\mu\omega(T_{\mu\nu})=0
\]
for all Hadamard states in the declared class.
\end{definition}

The trace has the form
\[
  \omega(T^\mu{}_\mu)
  =
  m^2\,\omega(:\!\Phi^2\!:)
  +(\xi-\xi_c)\,\omega(\Box:\!\Phi^2\!:)
  +\mathcal A[g,m,\xi],
\]
where \(\xi_c=(D-2)/(4D-4)\) and \(\mathcal A\) is a local curvature scalar
determined up to the same finite renormalization freedom.  In the massless
conformally coupled case, the state-dependent classical trace terms vanish
and the remaining local term is the trace anomaly.

\section{Relation to Flat-Space Composite Operators}

Point splitting is an operator-definition scheme.  It differs from defining a
regularized composite operator by inserting a classical functional into a
regulated path integral and then renormalizing.  The two schemes can be
compared by finite local counterterms when both are defined on the same
background and state class.  This comparison is the curved-background version
of the operator-coordinate changes developed in the renormalization volume.

\begin{openproblem}[Interacting stress tensors beyond formal perturbation]
\label{op:interacting-stress-tensors-beyond-formal-perturbation}
Construct stress tensors for interacting curved-spacetime QFTs with
nonperturbative control of domains, local covariance, conservation, trace
anomalies, and state dependence.  A solution should identify the exact
relation between locally covariant composite fields and flat-space
renormalized operator schemes.
\end{openproblem}
